\section{Limitations and Ethics}
\label{sec:limitations-ethics}

\emph{Scope of the analysis.}
All analyses use frozen encoders and a linear head, and the fine-tuning check covers one audio encoder with low-rank adapters on four targets over three iterations.
The two-cloud model assumes Gaussian score differences with a shared variance within each emotion pair, and \eqref{eq:lambda}, \eqref{eq:value}, and Proposition~\ref{prop:design} rest on first- or second-order approximations; we test their predictions empirically rather than claiming exact results.
The threshold-first account holds for single-modality audio and face features at $K \le 5$, not for fused and text configurations at $K=5$, and why the threshold share is large for the face, whose two losses are about equal, remains unexplained.

\emph{Data.}
We study four emotions in seven groups from six English corpora, which define their labels differently (Section~\ref{sec:setting-data}); speaker separation in CMU-MOSEI is incomplete, and RAVDESS has no face features in our pipeline.
Configurations with text have four groups, so the method comparisons with text rest on four targets, and the second encoders are tested for the method only in the AVT configuration.

\emph{Statistics and selection.}
Our intervals and tests reflect the draw of labels and test speakers within our corpora (Section~\ref{sec:setting-selection}).
Selecting hyperparameters on pseudo-targets loses up to \NUM{cmp-oursx-oracle-vs-tuned-avt-t-k0} points against the best value on the test speakers at $K=0$ (Section~\ref{sec:comparison-selection}).
Column normalization equalizes the per-emotion totals although the pools of our targets are skewed (Table~\ref{tab:data}, e.g., MELD); we did not analyze how its gain depends on the class ratios, and stronger skew can make it harmful~\cite{veilleux2021realistic}.

\emph{Ethics.}
We used only public corpora under their distribution terms and recorded no new data.
Faces enter the pipeline only as embeddings of detected face crops, and we release no identifying data.
Emotion recognition can be misused for surveillance or for decisions about people without their consent; the methods here estimate categorical labels from acted and media recordings and are not validated for decisions about individuals.
